\chapter{活动标架与 Gauss--Codazzi--Ricci 方程}

前面所有曲率都属于流形自身的内在几何。若一个 Riemann 流形 $M^m$ 浸入更高维 Riemann 流形 $\overline M^{m+q}$，环境 Levi--Civita 联络沿 $M$ 的导数既有切向分量，也有法向分量。后者就是外在弯曲。最一般的兼容条件不是只适用于曲面的 Gauss--Codazzi 两式，而是任意余维的 Gauss、Codazzi 与 Ricci 三个曲率块。本章先在一般 $(m,q)$ 情形推导它们，再把 $M^2\subset\RR^3$ 作为余维一、平坦环境的特例。

\section{切向与法向联络的分块}

设
\[
F:M^m\longrightarrow\overline M^{m+q}
\]
为等距浸入。以后把 $TM$ 与其像识别。沿 $M$ 有正交分解
\[
T\overline M|_M=TM\oplus NM,
\]
其中 $NM$ 是法丛。环境度量与 Levi--Civita 联络分别记为 $\overline g$、$\overline\nabla$，诱导度量与 Levi--Civita 联络记为 $g$、$\nabla$。

对切向量场 $X,Y\in\Gamma(TM)$，把环境导数分解为
\begin{equation*}{
\overline\nabla_XY
=\nabla_XY+B(X,Y),}
\end{equation*}
其中
\[
B(X,Y):=(\overline\nabla_XY)^\perp\in\Gamma(NM)
\]
称为第二基本形式。环境联络无挠，且 $\nabla$ 也是无挠的，因此
\[
B(X,Y)-B(Y,X)
=(\overline\nabla_XY-\overline\nabla_YX-[X,Y])^\perp=0,
\]
故 $B$ 对称。

对法向场 $\xi\in\Gamma(NM)$，环境导数同样分解为
\begin{equation*}{
\overline\nabla_X\xi
=-A_\xi X+\nabla_X^\perp\xi.}
\end{equation*}
这里 $A_\xi:TM\to TM$ 是关于法向 $\xi$ 的 形状算符，$\nabla^\perp$ 是 正规 联络。由 $X\overline g(Y,\xi)=0$ 与度量相容，
\[
0=\overline g(\overline\nabla_XY,\xi)
+\overline g(Y,\overline\nabla_X\xi),
\]
所以
\begin{equation}\label{eq:math-dg-shape-second-form-relation}
{
\overline g(B(X,Y),\xi)
=g(A_\xi X,Y).}
\end{equation}
由于 $B$ 对称，每个 $A_\xi$ 都关于 $g$ 自伴。

取适配正交标架
\[
\{e_i\}_{i=1}^m\subset TM,
\qquad
\{n_\alpha\}_{\alpha=1}^q\subset NM.
\]
写
\[
B(e_i,e_j)=h^\alpha{}_{ij}n_\alpha.
\]
则
\[
h^\alpha{}_{ij}=h^\alpha{}_{ji}.
\]
平均曲率向量定义为
\begin{equation*}\vb{H}
:=\frac1m\operatorname{tr}_gB.
\end{equation*}
它是法向量场，而不是内在标量曲率。




\begin{workedexample}
对单位球面取外法向 \(n=r\)，按本书约定 \(\overline\nabla_Xn=-AX\)，求形状算符、主曲率、高斯 曲率和有向平均曲率；再反转法向。

\textbf{解。} 因 \(D_Xn=X\)，故 \(A=-I\)。两个主曲率均为 \(-1\)，\(K=\det A=1\)，而 \(\mathsf H=\tfrac12\operatorname{tr}A=-1\)。反转法向后 \(A\mapsto-A\)，所以主曲率与有向平均曲率变号，\(K=\det A\) 不变。
\end{workedexample}

\section{Gauss--Codazzi--Ricci 的统一推导}

记 $\overline R$、$R$ 与 $R^\perp$ 分别为环境联络、切向 Levi--Civita 联络和 正规 联络 的曲率，均采用
\[
R(X,Y)Z
=\nabla_X\nabla_YZ-\nabla_Y\nabla_XZ-\nabla_{[X,Y]}Z
\]
的符号约定。对第二基本形式定义协变导数
\begin{equation*}(\nabla_X^\perp B)(Y,Z)
:=\nabla_X^\perp(B(Y,Z))
-B(\nabla_XY,Z)-B(Y,\nabla_XZ).
\end{equation*}

三条方程分别来自环境曲率的三个块：切--切块给 Gauss，切--法混合块给 Codazzi，法--法块给 Ricci。高余维时 形状算符 之间可能不交换，Ricci 方程正是 正规 丛 曲率与这种不交换性的联系。

\begin{workedexample}
从 Gauss--Weingarten 分解重新计算环境曲率的切向、法向部分，并核对一般余维 Gauss、Codazzi、Ricci 方程及交换子次序。

\textbf{解。} 将 \(\overline\nabla_YZ=\nabla_YZ+B(Y,Z)\) 代入 \(\overline\nabla_X\overline\nabla_YZ\)，再减去 \(X,Y\) 互换项及 \(\overline\nabla_{[X,Y]}Z\)。切向投影给出 Gauss 方程，法向投影给出 Codazzi 方程。对法向场 \(\xi\) 使用 \(\overline\nabla_Y\xi=-A_\xi Y+\nabla_Y^\perp\xi\)，取法向投影并与 \(\eta\) 内积，得到
\[
\langle R^\perp(X,Y)\xi,\eta\rangle
=\langle\overline R(X,Y)\xi,\eta\rangle+\langle[A_\xi,A_\eta]X,Y\rangle.
\]
故交换子次序为 \([A_\xi,A_\eta]\)，与本章末的完整分块计算一致。
\end{workedexample}


\section{移动标架形式}

在适配正交标架 $\{e_i,n_\alpha\}$ 中，环境联络一形式分块为
\[
\overline\omega
=
\begin{pmatrix}
\omega^i{}_j & \overline\omega^i{}_\beta\\
\overline\omega^\alpha{}_j & \omega^\alpha{}_\beta
\end{pmatrix}.
\]
其中 $\omega^i{}_j$ 是 $TM$ 的 Levi--Civita 联络，$\omega^\alpha{}_\beta$ 是 正规 联络，而混合块编码第二基本形式：
\begin{equation}\label{eq:math-dg-moving-frame-second-form}
\overline\omega^\alpha{}_i
=h^\alpha{}_{ij}\theta^j,
\qquad
\overline\omega^i{}_\alpha
=-h^\alpha{}_{ij}\theta^j.
\end{equation}
把这一分块代入环境第二结构方程
\[
\overline\Omega
=\dd\overline\omega+\overline\omega\wedge\overline\omega
\]
并分别读取切--切、切--法与法--法块，就得到\eqnrefs{eq:math-dg-gauss-general}、\eqnrefs{eq:math-dg-codazzi-general}与\eqnrefs{eq:math-dg-ricci-general}的外形式版本。移动标架法并不是另一套理论，而是同一分块曲率计算的紧凑表示。

上述三条兼容方程不仅是任何浸入都必须满足的必要条件；在平坦环境中，它们在适当拓扑假设下也是重构浸入的充分条件。这一事实称为子流形基本定理，它说明 Gauss--Codazzi--Ricci 并不是若干互不相关的恒等式，而是移动标架系统的完整可积条件。

设 $(M^m,g)$ 为单连通 Riemann 流形，$N\to M$ 是秩 $q$ 的 欧氏 向量丛，带 度量 联络 $\nabla^\perp$，并给定对称双线性形式
\[
B\in\Gamma\bigl(\operatorname{Sym}^2T^*M\otimes N\bigr).
\]
若 $(g,\nabla^\perp,B)$ 满足 欧氏 环境 $\overline R=0$ 时的 Gauss、Codazzi、Ricci 方程，则局部存在等距浸入
\[
F:M\longrightarrow\RR^{m+q}
\]
实现这些数据；在单连通情形可作全局重构，并且在固定初始点、初始适配标架以后唯一。去掉初始归一化后，唯一性只差一个 欧氏 刚体运动。

子流形的 Gauss--Codazzi--Ricci 三方程由此等价于扩展联络的平坦性 $\mathcal F=\dd\mathcal A+\mathcal A\wedge\mathcal A=0$（\eqnrefs{eq:math-dg-submanifold-flat-extended-connection}）。

\section{欧氏 三维空间中的曲面}

现在令 $m=2,q=1$，并取环境 $\overline M=\RR^3$ 的 欧氏 度量。设 $(e_1,e_2,n)$ 为适配正交归一标架，余标架为 $(\theta^1,\theta^2,\theta^3)$。限制到曲面后
\[
\theta^3=0.
\]
环境第一结构方程的法向分量给出
\begin{equation*}\omega^3{}_1\wedge\theta^1
+\omega^3{}_2\wedge\theta^2=0.
\end{equation*}
写
\[
\omega^3{}_1=h_{11}\theta^1+h_{12}\theta^2,
\qquad
\omega^3{}_2=h_{21}\theta^1+h_{22}\theta^2,
\]
代回上式得到
\begin{equation*}h_{12}=h_{21},
\end{equation*}
这正是一般结论“$B$ 对称”在二维余维一移动标架中的分量形式。第二基本形式写成
\begin{equation*}II=h_{ij}\theta^i\theta^j.
\end{equation*}
按本书约定，形状算符 $A$ 满足
\[
\overline\nabla_Xn=-AX,
\qquad
II(X,Y)=g(AX,Y).
\]
Gauss--Weingarten 公式于是成为
\begin{equation*}{
\begin{aligned}
\overline\nabla_XY&=\nabla_XY+II(X,Y)n,\\
\overline\nabla_Xn&=-AX.
\end{aligned}}
\end{equation*}

因为 欧氏 环境平坦，$\overline R=0$。Gauss 方程在正交基 $e_1,e_2$ 上给出
\[
K
=g(R(e_1,e_2)e_2,e_1)
=h_{11}h_{22}-h_{12}^2,
\]
即
\begin{equation*}{K=\det A.}
\end{equation*}
若 $\kappa_1,\kappa_2$ 是主曲率，则
\[
K=\kappa_1\kappa_2,
\qquad
\mathsf H=\frac12(\kappa_1+\kappa_2),
\]
其中 $\mathsf H$ 表示有向标量平均曲率，平均曲率向量为 $\vb{H}=\mathsf H n$。采用不同法向会使 $\kappa_i$ 与 $\mathsf H$ 同时变号，但 $K$ 不变。

Codazzi 方程在平坦环境中化为
\[
(\nabla_X^\perp B)(Y,Z)
=(\nabla_Y^\perp B)(X,Z).
\]
余维一时 $\nabla^\perp n=0$，若定义
\[
Dh_{ij}
:=\dd h_{ij}-h_{kj}\omega^k{}_i-h_{ik}\omega^k{}_j
=h_{ij;k}\theta^k,
\]
则移动标架的两个独立方程可写为
\begin{equation*}\dd\omega^3{}_1+\omega^3{}_2\wedge\omega^2{}_1=0,
\end{equation*}
\begin{equation*}\dd\omega^3{}_2+\omega^3{}_1\wedge\omega^1{}_2=0,
\end{equation*}
并等价于分量式
\begin{equation*}{h_{ij;k}=h_{ik;j}.}
\end{equation*}
余维一只有一个 形状算符，所以 Ricci 方程中的交换子自动为零；正规 丛 也是一维，Ricci 方程不再给出额外独立条件。

对定向 欧氏 超曲面
\[
F:M^m\to\RR^{m+1},
\]
单位法向本身定义 Gauss 映射
\begin{equation*}\mathcal N:M\longrightarrow S^m,
\qquad
\mathcal N(p)=n(p).\end{equation*}
由 Weingarten 公式，任意 $X\in T_pM$ 都有
\begin{equation}
{
\dd\mathcal N_p(X)
=\overline\nabla_Xn
=-A_pX.
}
\label{eq:math-dg-gauss-map-differential}
\end{equation}
所以 形状算符 不是抽象定义出来的矩阵：它就是 Gauss 映射 的负微分。其本征方向是法向变化最快的正交方向，主曲率则是这些方向上的有符号伸缩率。

把单位球面的标准度量记为 $g_{S^m}$，则 Gauss 映射 拉回度量为
\begin{equation*}(\mathcal N^*g_{S^m})(X,Y)
=g(AX,AY).\end{equation*}
右边常称为第三基本形式。另一方面，Gauss 映射 的有向 Jacobian 为
\begin{equation*}\det(\dd\mathcal N)
=(-1)^m\det A.\end{equation*}
对曲面 $M^2\subset\RR^3$，$\det A=K$，故 高斯 曲率 恰是 Gauss 映射 的局部面积伸缩率。$K$ 的“内在性”和法向映射的“外在定义”在这里并不矛盾：Gauss 方程保证二者数值一致。

形状算符 还精确控制平行超曲面的首次奇性。定义
\[
F_s(p):=F(p)+s\,n(p).
\]
对 $X\in T_pM$，由\eqnrefs{eq:math-dg-gauss-map-differential}
\begin{equation*}\dd F_s(X)
=(I-sA)X.\end{equation*}
因此只要 $I-sA$ 可逆，$F_s$ 仍是局部浸入；在主方向 $e_i$ 上，
\[
\dd F_s(e_i)=(1-s\kappa_i)e_i.
\]
当 $s=1/\kappa_i$ 时该方向退化，这就是 焦点 出现的外在版本。

在未退化区域内，法向仍可选为同一个 $n$。由
\[
-A_s\bigl((I-sA)X\bigr)
=\overline\nabla_Xn
=-AX
\]
得到
\begin{equation*}{
A_s=A(I-sA)^{-1},
\qquad
\kappa_i(s)=\frac{\kappa_i}{1-s\kappa_i}.
}\end{equation*}
同时面积元按
\begin{equation}
\operatorname{vol}_{g_s}
=|\det(I-sA)|\,\operatorname{vol}_g
=\prod_{i=1}^m|1-s\kappa_i|\,\operatorname{vol}_g
\label{eq:math-dg-parallel-area-jacobian}
\end{equation}
变化。展开这个乘积就会依次出现平均曲率、第二基本形式的初等对称多项式以及 Gauss--Kronecker 曲率。于是 tube 公式 与曲率积分的系数，本质上都来自同一个 行列式。

\begin{proposition}[平行浸入的二阶面积响应]
在无焦点的平行浸入范围，记 $\sigma_k$ 为主曲率的第 $k$ 个初等对称多项式。
则
\[
\begin{aligned}
\det(I-sA)
&=1-s\sigma_1+s^2\sigma_2+O(s^3)\\
&=1-sm\mathsf H+\frac{s^2}{2}
\bigl(m^2\mathsf H^2-|A|^2\bigr)+O(s^3).
\end{aligned}
\]
\end{proposition}

\begin{proof}
展开 $\prod_i(1-s\kappa_i)$，二阶系数为
$\sigma_2=\sum_{i<j}\kappa_i\kappa_j$。
由
\[
\sigma_1^2=\sum_i\kappa_i^2+2\sum_{i<j}\kappa_i\kappa_j
=|A|^2+2\sigma_2
\]
得到第二行。不能再写成 $\tfrac12(\sigma_1^2-\sigma_2)$，
因为 $\sigma_2$ 是初等对称函数，不是平方范数。
\end{proof}

\begin{note}[平行面积、管体积与稳定性]
二阶系数消失只说明 $\sigma_2=0$，不能断言全测地；
圆柱的主曲率只有一个非零，就给出直接反例。
欧氏 紧子流形的嵌入管状邻域在足够小的 $r$ 下满足
\[
\operatorname{Vol}(T_r(M))
=\int_M\int_{\{v\in N_pM:\ |v|<r\}}
\det(I-A_v)\,\dd v\,\operatorname{vol}_g(p).
\]
法向球关于 $v\mapsto-v$ 对称，所以法向积分中的奇次项消失，
首项为 $\omega_{n-m}r^{n-m}\operatorname{Vol}(M)$。
这是 Weyl 管公式的局部 Jacobian 依据；高余维不能用单一法向的
$\prod_i(1-s\kappa_i)$ 和一维 $\dd s$ 代替完整法向球积分。
管公式与热核系数都包含曲率不变量，但组合及归一化须分别计算。

极小条件使一阶面积变化消失。稳定性必须检验合法紧支或固定边界的
法向速度 $V=\phi n$，其二阶变分还含正项 $|\nabla\phi|^2$；
只看平行推进的 $-|A|^2$ 系数，不能断言多数肥皂膜不稳定。
边界约束与完整 Jacobi 二次型共同决定 Plateau 解的稳定性。
\end{note}

\begin{exbox}[球面与圆柱的外法向平行推进]
取外法向和约定 $A=-\dd n$，求球面、圆柱面沿法向平行推进后的面积元，核对一阶变化与平均曲率的关系。

\textbf{解。} 取外法向且 $A=-\dd n$。半径 $R$ 的球面主曲率均为 $-1/R$，因此
\[
\operatorname{vol}_{g_s}=(1+s/R)^m\operatorname{vol}_g,
\qquad\mathsf H=-1/R.
\]
其一阶系数 $m/R=-m\mathsf H$ 与面积变分吻合。
圆柱 $S^1(R)\times\RR\subset\RR^3$ 的主曲率为 $(-1/R,0)$，
面积比为 $1+s/R$，一阶系数为 $-2\mathsf H=1/R$，二阶项为零。
\end{exbox}

外在几何真正进入变分问题时，第二基本形式不再只是兼容条件中的系数，而直接控制面积的一阶与二阶变化。设
\[
F_t:M^m\to\overline M^{m+q}
\]
是一族浸入，变分速度为
\[
V:=\partial_tF_t|_{t=0}.
\]
把 $V=V^\top+V^\perp$ 分解成切向与法向部分。切向部分只对应参数重标记；内在面积真正改变的主项来自 $V^\perp$。对紧支撑变分，一阶体积公式为
\begin{equation}
{
\left.\frac{\dd}{\dd t}\operatorname{Vol}(F_t(M))\right|_{t=0}
=-m\int_M
\overline g(\vb{H},V^\perp)\,\operatorname{vol}_g
}.
\label{eq:math-dg-first-variation-submanifold-volume}
\end{equation}
若有边界，还要加上由 $V^\top$ 产生的边界通量项。因此
\[
{\vb{H}=0}
\]
恰是对所有紧支撑法向变分面积临界的 Euler--Lagrange 方程；这样的浸入称为 极小子流形。

对余维一 极小 超曲面，取单位法向 $n$ 与法向变分
\[
V=\phi n.
\]
在 $\vb{H}=0$ 的临界点，一阶变分消失，二阶变分由 稳定性 二次型 控制：
\begin{equation}
{
\delta^2\operatorname{Vol}[\phi]
=
\int_M
\left(
|\nabla\phi|^2
-\bigl(|A|^2+\overline{\operatorname{Ric}}(n,n)\bigr)\phi^2
\right)\operatorname{vol}_g
}.
\label{eq:math-dg-minimal-second-variation}
\end{equation}
这与前面测地线的 指标 形式 具有完全相同的结构：梯度项提供正的``弹性''，曲率势项可能制造负方向。事实上，一维 极小子流形 就是测地线，而\eqnrefs{eq:math-dg-minimal-second-variation}正退化为测地线 指标 形式。

若\eqnrefs{eq:math-dg-minimal-second-variation}对所有紧支撑 $\phi$ 都非负，称 极小 超曲面 稳定。因而稳定性问题被化成一个 Schrödinger 型二次型的谱下界问题；前面 Hilbert 空间、自伴算子与 Friedrichs 实现的理论在这里重新出现。几何变分、谱理论和 PDE 并不是三块独立知识，而是在这一 Jacobi 二次型 上闭合。

把分部积分显式做出来，二阶变分可写成
\begin{equation*}\delta^2\operatorname{Vol}[\phi]
=\int_M\phi\,L_J\phi\,\operatorname{vol}_g,
\qquad
L_J:=-\Delta_{\mathrm{LB}}
-\left(|A|^2+\overline{\operatorname{Ric}}(n,n)\right).\end{equation*}
这里 $L_J$ 是与当前 Laplace--Beltrami 符号约定相容的 自伴 Jacobi 算符。稳定性等价于 $L_J$ 的 二次型 非负；Morse 指标 则等于其负本征值总数（计重数）。这与测地线的 Morse 指标定理 完全平行，只是普通 Sturm--Liouville 算子被 超曲面 上的 Schrödinger 算符 取代。

欧氏 极小 超曲面 还有一个重要的二阶恒等式，把第二基本形式本身变成椭圆方程的未知量。若 $\overline M=\RR^{m+1}$ 且 $\mathsf H=0$，则 Codazzi 与 Gauss 方程推出 Simons 恒等式
\begin{equation}
{
\Delta_{\mathrm{LB}}h_{ij}
+|A|^2h_{ij}=0.
}
\label{eq:math-dg-simons-hessian}
\end{equation}
与 $h_{ij}$ 收缩，并使用
\[
\frac12\Delta_{\mathrm{LB}}|A|^2
=\langle\Delta_{\mathrm{LB}}A,A\rangle
+|\nabla A|^2,
\]
得到标量形式
\begin{equation*}{
\frac12\Delta_{\mathrm{LB}}|A|^2
=|\nabla A|^2-|A|^4.
}\end{equation*}
它说明 极小 条件并没有让外在曲率变成 调和；$|A|^4$ 是一个临界的非线性反应项，恰好与梯度耗散竞争。

\begin{note}[Simons 恒等式的适用环境与 Bernstein 定理]
Bernstein 定理说：$\RR^m$ 上的整 欧氏 极小图在 $m\le7$ 时必为超平面，
$m\ge8$ 已有非线性整极小图反例。
这个维数门槛来自稳定性、极小锥及尺度分析，不能由
$\tfrac2m-1$ 的符号推出；该系数早在 $m=2$ 就改变符号。
仅凭 $|A|^2$ 有界和上面的标量恒等式，也不能宣称已经证明整图平坦。
Simons 恒等式提供局部曲率估计的输入，完整全局结论还需其他分析步骤。

若环境是单位球面，恒等式会增加环境曲率项：极小超曲面满足
\[
\tfrac12\Delta_{\mathrm{LB}}|A|^2
=|\nabla A|^2+m|A|^2-|A|^4.
\]
因此 $S^3$ 中的 Clifford 极小环面有 $m=2$、$|A|^2=2$、$\nabla A=0$，
右端为 $0+4-4=0$。它不能用没有 $m|A|^2$ 项的 欧氏 公式核验；
当作 $\RR^4$ 中余维二的环面时，其 欧氏 平均曲率也并非零。

Chern--Simons 形式
$\operatorname{CS}(A)=\operatorname{tr}(A\wedge\dd A+\tfrac23A\wedge A\wedge A)$
满足外微分恒等式 $\dd\operatorname{CS}(A)=\operatorname{tr}(F\wedge F)$；
其一阶变分则为 $2\operatorname{tr}(\delta A\wedge F)$ 加全导数，
不能把外微分与泛函变分混为一谈。
在 欧氏 世界面的面积模型中，极小方程与 Jacobi 算符也控制一阶、二阶几何变分；
这可作为弦世界面分析的几何输入，量子修正仍需独立推导。
\end{note}

对图形浸入，极小 方程可以直接写成一个拟线性椭圆 PDE。令
\[
F(x)=(x,u(x)),
\qquad x\in\Omega\subset\RR^m.
\]
诱导度量满足
\[
g_{ij}=\delta_{ij}+u_i u_j,
\qquad
\det g=1+|\nabla u|^2.
\]
因此面积泛函为
\begin{equation*}\mathcal A[u]
=\int_\Omega\sqrt{1+|\nabla u|^2}\,\dd^m x.\end{equation*}
做紧支撑变分 $u\mapsto u+\varepsilon\eta$，有
\[
\begin{aligned}
\delta\mathcal A[u;\eta]
&=\int_\Omega
\frac{\nabla u\cdot\nabla\eta}
{\sqrt{1+|\nabla u|^2}}\,\dd^mx\\
&=-\int_\Omega
\eta\,
\operatorname{div}\left(
\frac{\nabla u}{\sqrt{1+|\nabla u|^2}}
\right)\dd^mx.
\end{aligned}
\]
所以 极小 图 方程 是
\begin{equation*}{
\operatorname{div}\left(
\frac{\nabla u}{\sqrt{1+|\nabla u|^2}}
\right)=0
}.\end{equation*}
在小斜率 $|\nabla u|\ll1$ 下展开
\[
\frac{1}{\sqrt{1+|\nabla u|^2}}
=1-\frac12|\nabla u|^2+O(|\nabla u|^4),
\]
最低阶退化为
\[
\Delta u=0.
\]
因此 调和 图 是 极小 图 的线性化，而不是与其无关的偶然近似。

面积最速下降则给出 平均曲率 流。若选择纯法向速度
\[
V=m\vb{H},
\]
则由一阶变分公式
\[
\frac{\dd}{\dd t}\operatorname{Vol}(M_t)
=-m^2\int_{M_t}|\vb{H}|^2\,\operatorname{vol}_{g_t}\le0.
\]
这说明 平均曲率 流 与 Ricci 流 在逻辑上不同：前者演化的是浸入并以外在平均曲率降低面积，后者直接演化内在度量。两者都可看作几何泛函的梯度型演化，但状态空间和 规范 结构不同。

\begin{exbox}[单位球面与圆柱面]
按 $A=-\dd n$ 的约定，求单位球面与半径 $R$ 圆柱面的主曲率、高斯曲率和平均曲率。

\textbf{解。} 单位球面取外法向 $n=r$。对切向量 $X$，
\[
\overline\nabla_Xn=X.
\]
由约定 $\overline\nabla_Xn=-AX$，有 $A=-I$，故
\[
\kappa_1=\kappa_2=-1,
\qquad
K=1,
\qquad
\mathsf H=-1.
\]

半径 $R$ 的圆柱面取外法向时，圆周方向主曲率为 $-1/R$，轴向主曲率为零，因此
\[
K=0,
\qquad
\mathsf H=-\frac{1}{2R}.
\]
圆柱具有非零外在弯曲却有零 高斯 曲率，直接说明平均曲率依赖嵌入，而 $K$ 是内在度量可以检测的曲率。
\end{exbox}

一般余维的三条兼容方程到此才在 $q=1$、$\overline R=0$ 的条件下退化为熟悉的曲面公式。这样的顺序保证 Gauss--Codazzi 不是从三维曲面的特殊计算“推广”出来，而是环境联络曲率分块的一般结论。
Gauss--Codazzi--Ricci 方程已经分离切向和法向曲率。曲面的显式计算把内禀曲率、第二基本形式和环境几何重新联系起来。

\section{曲面嵌入后：内禀曲率与外在弯曲}

球面内部可以测距离和曲率；作为 \(\RR^3\) 中的曲面，还可以观察法向怎样变化。前者只用诱导度量，后者必须用环境导数。两者并非独立。

\begin{definition}[第二基本形式与法向联络]
设 \(M\) 是 Riemann 流形 \((\overline M,\overline g)\) 的嵌入子流形，\(\overline\nabla\) 为环境 Levi--Civita 联络。环境切空间正交分解为 \(T_p\overline M=T_pM\oplus(T_pM)^\perp\)。对切向量场 \(X,Y\)，定义
\[
\overline\nabla_XY=\nabla_XY+\mathrm{II}(X,Y),
\]
其中两项分别是切向、法向投影。法向场 \(\xi\) 的法向导数为
\(\nabla_X^\perp\xi=(\overline\nabla_X\xi)^\perp\)。
\end{definition}

\(\mathrm{II}\) 是双线性张量：\((\overline\nabla_X(fY))^\perp=f\mathrm{II}(X,Y)\)，因为额外项 \((Xf)Y\) 为切向。环境联络无挠还给出 \(\mathrm{II}(X,Y)=\mathrm{II}(Y,X)\)。与此不同，\(\nabla_X(fY)\) 保留 \((Xf)Y\)，所以联络本身不是张量。

定义形算子 \(A_\xi:TM\to TM\) 使
\(\langle A_\xi X,Y\rangle=\langle\mathrm{II}(X,Y),\xi\rangle\)。
对 \(\langle\xi,Y\rangle=0\) 求导，利用度量相容性，便得到 Weingarten 分解
\[
\overline\nabla_X\xi=-A_\xi X+\nabla_X^\perp\xi.
\]

采用 \(R(X,Y)Z=\nabla_X\nabla_YZ-\nabla_Y\nabla_XZ-\nabla_{[X,Y]}Z\) 的曲率约定。把 \(\overline R(X,Y)Z\) 中的每一次环境导数按上述两式分解。先把一项算到底：由切法分解与 Weingarten 分解各用一次，
\[
\overline\nabla_X\overline\nabla_YZ
=\overline\nabla_X\bigl(\nabla_YZ+\mathrm{II}(Y,Z)\bigr)
=\underbrace{\nabla_X\nabla_YZ}_{\text{切向}}
+\underbrace{\mathrm{II}(X,\nabla_YZ)}_{\text{法向}}
\underbrace{-A_{\mathrm{II}(Y,Z)}X}_{\text{切向}}
+\underbrace{\nabla_X^\perp\mathrm{II}(Y,Z)}_{\text{法向}}.
\]
交换 \(X,Y\) 再相减，并代入
\(\overline\nabla_{[X,Y]}Z=\nabla_{[X,Y]}Z+\mathrm{II}([X,Y],Z)\)。
收集切向部分：两个 \(\nabla\nabla Z\) 项与 \(\nabla_{[X,Y]}Z\)
合成 \(R(X,Y)Z\)，两个 Weingarten 项幸存，得到
\[
(\overline R(X,Y)Z)^\top
=R(X,Y)Z-A_{\mathrm{II}(Y,Z)}X+A_{\mathrm{II}(X,Z)}Y.
\]
与 \(W\in TM\) 内积时，形算子的定义把 Weingarten 项化为
\(-\langle\mathrm{II}(Y,Z),\mathrm{II}(X,W)\rangle\)；
移项后便产生 Gauss 方程
\begin{equation}
\label{eq:math-gauss-compatibility}
\langle R(X,Y)Z,W\rangle
=\langle\overline R(X,Y)Z,W\rangle
+\langle\mathrm{II}(Y,Z),\mathrm{II}(X,W)\rangle
-\langle\mathrm{II}(X,Z),\mathrm{II}(Y,W)\rangle .
\end{equation}
同一展开的法向部分逐项为
\[
\begin{aligned}
(\overline R(X,Y)Z)^\perp
={}&\nabla_X^\perp\mathrm{II}(Y,Z)-\nabla_Y^\perp\mathrm{II}(X,Z)\\
&+\mathrm{II}(X,\nabla_YZ)-\mathrm{II}(Y,\nabla_XZ)
-\mathrm{II}([X,Y],Z).
\end{aligned}
\]
利用 \([X,Y]=\nabla_XY-\nabla_YX\)，便得到 Codazzi 方程
\[
(\overline R(X,Y)Z)^\perp
=(\nabla_X^\perp\mathrm{II})(Y,Z)
 -(\nabla_Y^\perp\mathrm{II})(X,Z),
\]
其中
\((\nabla_X^\perp\mathrm{II})(Y,Z)=\nabla_X^\perp(\mathrm{II}(Y,Z))
-\mathrm{II}(\nabla_XY,Z)-\mathrm{II}(Y,\nabla_XZ)\)。
后两项扣除了求导时切向基向量的变化。

对法向场 \(\xi\) 做同样的二次求导。先展开一项：
\[
\overline\nabla_X\overline\nabla_Y\xi
=\overline\nabla_X\bigl(-A_\xi Y+\nabla_Y^\perp\xi\bigr)
=-\nabla_X(A_\xi Y)-\mathrm{II}(X,A_\xi Y)
-A_{\nabla_Y^\perp\xi}X+\nabla_X^\perp\nabla_Y^\perp\xi ,
\]
其中第二项来自 Weingarten 分解中的切向项
\(-A_\xi Y\) 再求导：它的法向投影正是
\(-\mathrm{II}(X,A_\xi Y)\)；第三项来自对法向场
\(\nabla_Y^\perp\xi\) 再用一次 Weingarten 分解。
交换 \(X,Y\) 相减，并代入
\(\overline\nabla_{[X,Y]}\xi=-A_\xi[X,Y]+\nabla^\perp_{[X,Y]}\xi\)；
取法向部分后，三个二阶法向导数项合成
\(R^\perp(X,Y)\xi\)，切向的 \(-A_{\nabla_Y^\perp\xi}X\)
一类项不进入法向方程。于是得到
\[
(\overline R(X,Y)\xi)^\perp
=R^\perp(X,Y)\xi-\mathrm{II}(X,A_\xi Y)
+\mathrm{II}(Y,A_\xi X).
\]
与第二个法向场 \(\eta\) 内积，利用形算子的定义并移项，得到 Ricci 方程
\[
\langle R^\perp(X,Y)\xi,\eta\rangle
=\langle\overline R(X,Y)\xi,\eta\rangle
+\langle[A_\xi,A_\eta]X,Y\rangle .
\]
\(R^\perp\) 是法向联络的曲率。三式分别是同一次环境曲率计算的切向、混合与法向部分。

以半径 \(a>0\) 的球面为例，取外法向 \(n(p)=p/a\)。由于
\(\overline\nabla_Xn=X/a\)，有 \(A_n=-I/a\)，
\(\mathrm{II}(X,Y)=-\langle X,Y\rangle n/a\)。环境平直，在正交单位切向量 \(X,Y\) 上取 \(Z=Y,W=X\)，\cref{eq:math-gauss-compatibility} 给 \(K=1/a^2\)。法向反向使两个主曲率同时变号，乘积和内禀曲率保持不变。圆柱则只有一个非零主曲率，因此 \(K=0\)：外在弯曲不等于内禀曲率。

若曲面是函数图像，还能把这些抽象量化为通常的偏导数。设
\(F(x,y)=(x,y,f(x,y))\)，其中 \(f\) 光滑。切向量
\(F_x=(1,0,f_x)\)、\(F_y=(0,1,f_y)\) 给诱导度量
\[
(g_{ij})=\begin{pmatrix}1+f_x^2&f_xf_y\\f_xf_y&1+f_y^2\end{pmatrix},
\qquad \det g=1+f_x^2+f_y^2.
\]
取向上的单位法向量
\(n=(-f_x,-f_y,1)/\sqrt{1+f_x^2+f_y^2}\)。环境是平直的，二阶导数
\(F_{ij}=(0,0,f_{ij})\) 的法向分量为
\(\langle F_{ij},n\rangle=f_{ij}/\sqrt{1+|\nabla f|^2}\)。注意本章的约定是
\(\mathrm{II}(F_i,F_j)=(\overline\nabla_{F_i}F_j)^\perp\)，故
\(\mathrm{II}_{ij}=f_{ij}n/\sqrt{1+|\nabla f|^2}\)；对球面图像用向上法向时，这与先前外法向的符号可能相反，但曲率不受影响。

在二维，Gauss 方程说曲率等于形算子的行列式。确切地，对 \(X=F_x,Y=F_y,Z=F_y,W=F_x\) 代入 \cref{eq:math-gauss-compatibility}，环境曲率为零，得到
\[
\langle R(F_x,F_y)F_y,F_x\rangle
=\langle\mathrm{II}(F_y,F_y),\mathrm{II}(F_x,F_x)\rangle
-\|\mathrm{II}(F_x,F_y)\|^2.
\]
左边等于 \(K\det g\)，右边等于
\((f_{xx}f_{yy}-f_{xy}^2)/(1+|\nabla f|^2)\)。因此
\begin{equation}
\label{eq:math-graph-gaussian-curvature}
K=\frac{f_{xx}f_{yy}-f_{xy}^2}{(1+f_x^2+f_y^2)^2}.
\end{equation}
式 \cref{eq:math-graph-gaussian-curvature} 的分子是黑塞矩阵的行列式而非迹：两方向都向同一侧弯时它为正，鞍面一正一负时它为负。取 \(f(x,y)=\tfrac12(x^2+y^2)\)，原点给 \(K=1\)；取 \(f(x,y)=\tfrac12(x^2-y^2)\)，原点给 \(K=-1\)。这两个图像在原点有同一切平面，差别完全由二阶变化决定。

球面的计算还可逐项检查相容条件。参数化
\(F(\vartheta,\varphi)=a(\sin\vartheta\cos\varphi,
\sin\vartheta\sin\varphi,\cos\vartheta)\) 给诱导度量
\(g=a^2(\dd\vartheta^2+\sin^2\vartheta\,\dd\varphi^2)\)。
前章已由这个度量算得 \(R^\vartheta{}_{\varphi\vartheta\varphi}
=\sin^2\vartheta\)，故
\(K=R_{\vartheta\varphi\varphi\vartheta}/\det g=1/a^2\)，
与嵌入计算吻合。这里
\(\mathrm{II}=-g\,n/a\) 且 \(\nabla g=0\)；
法向 \(n\) 的法向导数为零，所以
\(\nabla^\perp\mathrm{II}=0\)。环境曲率也为零，
Codazzi 方程两侧均为零。

Ricci 方程在余维一曲面上较难看出内容，因为法空间只有一个方向，
两个形算子的交换子自动为零。在余维二可以做一个点上的计算。
令 \(F(u,v)=(u,v,(u^2-v^2)/2,uv)\) 把平面嵌入
\(\RR^4\)。在原点，切基为 \(e_1,e_2\)，法基为
\(n_1=e_3,n_2=e_4\)；二阶偏导给
\[
A_{n_1}=
\begin{pmatrix}1&0\\0&-1\end{pmatrix},
\qquad
A_{n_2}=
\begin{pmatrix}0&1\\1&0\end{pmatrix}.
\]
于是 Gauss 方程给该点的
\(K=\det A_{n_1}+\det A_{n_2}=-2\)。
两个形算子不对易：
\([A_{n_1},A_{n_2}]=
\begin{psmallmatrix}0&2\\-2&0\end{psmallmatrix}\)，
故 Ricci 方程给
\(\langle R^\perp(e_1,e_2)n_1,n_2\rangle=-2\)。
这不是凭空多出一个曲率：它记录两个不同法方向的弯曲作用顺序
不能交换。



\par
\begin{workedexample}
在圆柱面 \((a\cos\theta,a\sin\theta,z)\) 上求两个主曲率，再用 Gauss 方程计算 \(K\)。将结果与把圆柱沿母线剪开并铺平相比较。
\textbf{解。} 取单位法向 \(n=(\cos\theta,\sin\theta,0)\)。参数切向量为 \(F_\theta=(-a\sin\theta,a\cos\theta,0)\)、\(F_z=(0,0,1)\)，第一基本形式系数 \(E=a^2,G=1,F=0\)。按 \(h_{ij}=\langle F_{ij},n\rangle\)，第二基本形式的圆周分量为 \(-a\)，其余分量为零，故主曲率为 \(-1/a,0\)；反转法向只改变符号，不改变 \(K=0\)。圆柱沿母线剪开可等距铺成平面，与零内禀曲率一致。
\end{workedexample}
\par


\subsection{叶状超曲面的法向与翘曲截面}

对余维一的叶片，前面的形状算子可直接与适配度规中的法向导数比较。取 $K_{ab}=\langle\nabla_{\partial_a}n,\partial_b\rangle$，则 $K_{ab}=-\langle B_{ab},n\rangle$；单位外法向球面的 $K$ 为正，而本章 $B$ 的法向分量为负。两种写法由这一等式转换。

\begin{figure}[htbp]
\centering
% Local TikZ geometry; styles inherited exclusively from the shared preamble.
\begin{tikzfig}[scale=.91]
\draw[gray!55] (-3.2,-.7) .. controls (-1.9,-1.25) and (1.5,-1.28) .. (3.2,-.7);
\draw[gray!65] (-3.1,1.0) .. controls (-1.4,.5) and (1.9,.55) .. (3.1,1.0);
\draw[thick] (0,-.7) ellipse (2.3 and .55);
\draw[thick] (0,1.05) ellipse (2.75 and .53);
\draw[gray!60,thin] (-2.3,-.7)--(-2.75,1.05) (2.3,-.7)--(2.75,1.05);
\draw[lecturearrow] (0,-.66)--(.15,.8) node[midway,right] {$n$};
\draw[lecturearrow] (.15,.8)--(1.25,.89) node[midway,above] {$N^a$};
\draw[lecturearrow] (-1.1,-.72)--(-1.1,.5) node[midway,left] {$N$};
\node[lecturelabel] at (2.7,1.85) {$\Sigma_{t+\dd t}$};
\node[lecturelabel] at (2.43,-1.56) {$\Sigma_t$};
\node[lecturelabel] at (-2.1,2.0) {$M\simeq I\times\Sigma$};
\end{tikzfig}

\caption{叶状分解中的相邻截面、单位法向 $n$、法向进动 $N$ 与切向位移 $N^a$。图中的椭圆只代表局部截面；一般叶不必具有椭圆形状。}\label{fig:ghs-foliation-normals}
\end{figure}

\cref{fig:ghs-foliation-normals} 说明了法向传播与沿叶的重新标记属于不同的几何操作；其中的曲面形状只是示意，真正进入 拉普拉斯算子 的是诱导度规与外曲率。
把一般几何量代入球面坐标，可以直接检验这些定义。
若 $h_{ab}(r,\theta)=a(r)^2\gamma_{ab}(\theta)$ 且 $\gamma$ 与 $r$ 无关，则 $K_{ab}=(a'/a)h_{ab}$，$K=(d-1)a'/a$。
代入 \cref{eq:ghs-laplace-evolving-slices}，得到
\[
 \Delta_g=\partial_r^2+(d-1)\frac{a'}a\partial_r+a^{-2}\Delta_\gamma.
\]
取 $a(r)=r$ 且 $\gamma$ 为标准球度规，才恢复 $\RR^d$ 中的球坐标公式。

实际的空间未必有径向中心。即便可以在一点附近构造测地球，这也仅给局部的“径向”参数；是否能够将径向量子化推广到整个空间，还需要分析作用量、边界态与演化核。这些问题将在下一组数学工具建立后处理。



\section{浸入几何的公式联系}
以下取 Riemann 环境中的等距浸入，并使用本章第二基本形式与形算子的符号。平坦环境的曲面公式是一般兼容方程的余维一特例。
\begin{note}[公式之间的联系]
对等距浸入 \(M\hookrightarrow\overline M\)，
\begin{align*}
\overline\nabla_XY&=\nabla_XY+B(X,Y),\\
\overline\nabla_X\xi&=-A_\xi X+\nabla_X^\perp\xi,\\
\overline g(B(X,Y),\xi)&=g(A_\xi X,Y).
\end{align*}
三条兼容方程见\eqnrefs{eq:math-dg-gauss-general}、\eqnrefs{eq:math-dg-codazzi-general}与\eqnrefs{eq:math-dg-ricci-general}。曲面 \(M^2\subset\RR^3\) 的平坦余维一特例为
\[
K=\det A,
\qquad
h_{ij;k}=h_{ik;j}.
\]
\end{note}

\begin{theorem}[浸入的 Gauss--Codazzi--Ricci 方程]
等距浸入的第二基本形式 $B$、形状算符 $A$ 与法向联络满足 \eqnrefs{eq:math-dg-gauss-general}、\eqnrefs{eq:math-dg-codazzi-general}、\eqnrefs{eq:math-dg-ricci-general}。三式分别来自环境曲率的切向、混合与法向分块。
\end{theorem}

\begin{proof}
先让环境曲率作用在切向量 $Z$。由 Gauss 公式，
\[
\overline\nabla_YZ
=\nabla_YZ+B(Y,Z).
\]
再沿 $X$ 求环境协变导数，并分别对切向与法向部分使用 Gauss--Weingarten 分解：
\[
\begin{aligned}
\overline\nabla_X\overline\nabla_YZ
={}&\nabla_X\nabla_YZ+B(X,\nabla_YZ)\\
&-A_{B(Y,Z)}X
+\nabla_X^\perp B(Y,Z).
\end{aligned}
\]
交换 $X,Y$，再减去 $\overline\nabla_{[X,Y]}Z$。切向部分得到
\[
(\overline R(X,Y)Z)^\top
=R(X,Y)Z
-A_{B(Y,Z)}X+A_{B(X,Z)}Y.
\]
与切向量 $W$ 做内积并用\eqnrefs{eq:math-dg-shape-second-form-relation}，
\[
\begin{aligned}
\overline g(\overline R(X,Y)Z,W)
={}&g(R(X,Y)Z,W)\\
&-\overline g(B(X,W),B(Y,Z))\\
&+\overline g(B(Y,W),B(X,Z)).
\end{aligned}
\]
移项后得到 Gauss 方程
\begin{equation}\label{eq:math-dg-gauss-general}
{
\begin{aligned}
g(R(X,Y)Z,W)
={}&\overline g(\overline R(X,Y)Z,W)\\
&+\overline g(B(X,W),B(Y,Z))\\
&-\overline g(B(X,Z),B(Y,W)).
\end{aligned}}
\end{equation}

同一个展开的法向部分为
\[
(\overline R(X,Y)Z)^\perp
=(\nabla_X^\perp B)(Y,Z)
-(\nabla_Y^\perp B)(X,Z),
\]
即 Codazzi 方程
\begin{equation}\label{eq:math-dg-codazzi-general}
{
(\overline R(X,Y)Z)^\perp
=(\nabla_X^\perp B)(Y,Z)
-(\nabla_Y^\perp B)(X,Z).}
\end{equation}

最后让环境曲率作用在法向场 $\xi$。由 Weingarten 公式，
\[
\overline\nabla_Y\xi
=-A_\xi Y+\nabla_Y^\perp\xi.
\]
再次沿 $X$ 求导并取法向部分：
\[
\begin{aligned}
(\overline\nabla_X\overline\nabla_Y\xi)^\perp
=-B(X,A_\xi Y)
+\nabla_X^\perp\nabla_Y^\perp\xi.
\end{aligned}
\]
交换 $X,Y$ 并减去 $[X,Y]$ 项，得到
\[
(\overline R(X,Y)\xi)^\perp
=R^\perp(X,Y)\xi
-B(X,A_\xi Y)+B(Y,A_\xi X).
\]
与另一法向场 $\eta$ 做内积，利用 形状算符 自伴性，
\[
\begin{aligned}
\overline g(\overline R(X,Y)\xi,\eta)
={}&\overline g(R^\perp(X,Y)\xi,\eta)\\
&-g(A_\eta X,A_\xi Y)
+g(A_\eta Y,A_\xi X)\\
={}&\overline g(R^\perp(X,Y)\xi,\eta)
-g([A_\xi,A_\eta]X,Y).
\end{aligned}
\]
因此 Ricci 方程为
\begin{equation}\label{eq:math-dg-ricci-general}
{
\overline g(R^\perp(X,Y)\xi,\eta)
=\overline g(\overline R(X,Y)\xi,\eta)
+g([A_\xi,A_\eta]X,Y).}
\end{equation}
交换子的次序已经由上述展开固定；若把 $[A_\xi,A_\eta]$ 反写，右端必须同时变号。

退化自检：一维法丛（如 $\RR^3$ 中曲面）中 $[A_\xi,A_\eta]=0$，Ricci 方程 的法向部分自动满足；$\RR^4$ 中的标准 Clifford 环面具有可同时对角化的两个法向形状算符，故它们的交换子为零，法丛平坦；这由具体形状算符检验，而不是把它们笼统记作 $\pm J$。
\end{proof}

\begin{theorem}[欧氏 浸入的局部基本定理]
在可缩坐标邻域上给定度量 $g$、秩 $q$ 的带度量法丛及度量联络、对称 $B$。若满足平直环境的 Gauss--Codazzi--Ricci 方程，则存在等距浸入 $F:M\to\RR^{m+q}$ 实现这些数据，唯一性差一个 欧氏 刚体运动。
\end{theorem}

\begin{proof}
在局部正交标架下，把 Levi--Civita 联络、正规 联络 与第二基本形式组合成一个 $\mathfrak{so}(m+q)$-值一形式
\[
\mathcal A=
\begin{pmatrix}
\omega^i{}_j & -h^i{}_{\beta k}\theta^k\\
 h^\alpha{}_{jk}\theta^k & \omega^\alpha{}_{\beta}
\end{pmatrix}.
\]
这里上下混合块用环境度量识别，符号选择与\eqnrefs{eq:math-dg-moving-frame-second-form}一致。计算其曲率
\[
\mathcal F
:=\dd\mathcal A+\mathcal A\wedge\mathcal A.
\]
切--切块恰是 Gauss 方程的左减右，切--法块恰是 Codazzi 方程，法--法块恰是 Ricci 方程。因此
\begin{equation}
\text{Gauss+Codazzi+Ricci}
\quad\Longleftrightarrow\quad
\mathcal F=0.
\label{eq:math-dg-submanifold-flat-extended-connection}
\end{equation}
也就是说，全部兼容条件等价于这个扩展联络平坦。

在单连通区域上，平坦矩阵联络可积分为一个正交矩阵场 $Q:M\to O(m+q)$，使
\[
Q^{-1}\dd Q=\mathcal A
\]
（若采用相反的矩阵 约定，等价地写成 $\dd Q=Q\mathcal A$）。令 $Q_1,\ldots,Q_m$ 为前 $m$ 个列向量，并定义 $\RR^{m+q}$-值一形式
\[
\alpha:=\sum_{i=1}^m Q_i\,\theta^i.
\]
计算外微分：
\[
\dd\alpha
=\sum_i \dd Q_i\wedge\theta^i
+\sum_i Q_i\,\dd\theta^i.
\]
$\dd Q_i$ 的切向联络部分与无挠第一结构方程
\[
\dd\theta^i+\omega^i{}_j\wedge\theta^j=0
\]
恰好抵消；法向部分含
\[
h^\alpha{}_{ij}\theta^j\wedge\theta^i,
\]
又因 $h^\alpha{}_{ij}=h^\alpha{}_{ji}$ 对称而为零。故
\[
\dd\alpha=0.
\]
单连通性使闭的一形式 $\alpha$ 可全局积分，所以存在 $F:M\to\RR^{m+q}$ 满足
\[
\dd F=\alpha.
\]
因为 $Q_i$ 正交归一，
\[
\langle \dd F(X),\dd F(Y)\rangle
=\sum_i\theta^i(X)\theta^i(Y)=g(X,Y),
\]
故 $F$ 是等距浸入。再由 $\dd Q=Q\mathcal A$ 的混合块读出其第二基本形式与 正规 联络 正是预先给定的 $B$ 与 $\nabla^\perp$。这完成了存在性重构的核心计算。

若两个重构具有相同初始点与初始适配标架，它们满足同一一阶系统，因此由唯一性完全一致；改变初始点和初始正交标架只对应环境中的平移与正交变换。于是唯一性恰好差一个 欧氏 刚体运动。
\end{proof}

\begin{proposition}[平行曲面的平均曲率响应]
对 欧氏 定向超曲面，取 $F_s=F+sn$、$A=-\dd n$，并限制在 $1-s\kappa_i>0$ 的范围。则\[\operatorname{vol}_{g_s}=\det(I-sA)\operatorname{vol}_g,\qquad\left.\partial_s\operatorname{vol}_{g_s}\right|_{s=0}=-m\mathsf H\operatorname{vol}_g.\]
\end{proposition}

\begin{proof}
由\eqnrefs{eq:math-dg-parallel-area-jacobian}，平行曲面的面积元为
\begin{equation*}
\operatorname{vol}_{g_s}=|\det(I-sA)|\,\operatorname{vol}_g
=\prod_{i=1}^m|1-s\kappa_i|\,\operatorname{vol}_g,
\end{equation*}
其中 $\kappa_1,\ldots,\kappa_m$ 是 形状算符 $A$ 的本征值即主曲率。在 $s=0$ 附近假设尚未遇到 焦点，即 $1-s\kappa_i>0$ 对所有 $i$ 成立，可去掉绝对值，
\begin{equation*}
\operatorname{vol}_{g_s}=\det(I-sA)\,\operatorname{vol}_g
=\prod_{i=1}^m(1-s\kappa_i)\,\operatorname{vol}_g.
\end{equation*}
把行列式按特征多项式在 $s=0$ 处 Taylor 展开。由 $\det(I-sA)$ 是 $s$ 的 $m$ 次多项式，常数项为 $\det I=1$，一次项为 $-s\operatorname{tr}A$（特征多项式系数与初等对称多项式的关系，一次项系数恰为负迹），故
\begin{equation*}
\det(I-sA)=1-s\operatorname{tr}A+O(s^2).
\end{equation*}
等价地，对乘积形式逐项展开，
\begin{equation*}
\prod_{i=1}^m(1-s\kappa_i)
=1-s\sum_{i=1}^m\kappa_i+O(s^2)
=1-s\operatorname{tr}A+O(s^2),
\end{equation*}
其中用了 $\operatorname{tr}A=\sum_i\kappa_i$。由平均曲率定义 $\mathsf H=\frac1m\operatorname{tr}A=\frac1m\sum_i\kappa_i$，
\begin{equation*}
\operatorname{tr}A=m\mathsf H
\end{equation*}
（余维一时 $\vb{H}=\mathsf H n$）。代回面积元展开，
\begin{equation*}
\operatorname{vol}_{g_s}=(1-sm\mathsf H+O(s^2))\,\operatorname{vol}_g.
\end{equation*}
对 $s$ 求导并在 $s=0$ 取值，
\begin{align*}
\left.\frac{\dd}{\dd s}\operatorname{vol}_{g_s}\right|_{s=0}
&=\left.\frac{\dd}{\dd s}\bigl[(1-sm\mathsf H+O(s^2))\,\operatorname{vol}_g\bigr]\right|_{s=0}\\
&=\bigl[-m\mathsf H+O(s)\bigr]\Big|_{s=0}\operatorname{vol}_g\\
&=-m\mathsf H\,\operatorname{vol}_g.
\end{align*}
这正是后面\eqnrefs{eq:math-dg-first-variation-submanifold-volume}在常法向速度 $V=n$ 情形的局部版本。换言之，平均曲率既是 形状算符 的迹，也是平行法向推进时面积元的一阶响应。
\end{proof}

\begin{proposition}[浸入体积的一阶变分]
设 $F_t$ 为浸入变分、$V=\partial_tF$、$\vb{H}=m^{-1}\operatorname{tr}_gB$。则\[\partial_t\operatorname{vol}_{g_t}=\bigl(\operatorname{div}_MV^\top-m\overline g(\vb{H},V^\perp)\bigr)\operatorname{vol}_g.\]闭子流形或紧支法向变分下，体积临界条件恰为 $\vb{H}=0$。
\end{proposition}

\begin{proof}
取局部坐标 $x^i$，诱导度量为
\[
g_{ij}(t)
=\overline g(\partial_iF_t,\partial_jF_t).
\]
时间求导并用环境联络度量相容性：
\[
\dot g_{ij}
=\overline g(\overline\nabla_iV,e_j)
+\overline g(e_i,\overline\nabla_jV).
\]
因此
\[
\frac12\operatorname{tr}_g\dot g
=g^{ij}\overline g(\overline\nabla_iV,e_j).
\]
将 $V=V^\top+V^\perp$ 分解。切向部分给出
\[
g^{ij}g(\nabla_iV^\top,e_j)
=\operatorname{div}_M V^\top.
\]
对法向部分，利用
\[
0=e_i\,\overline g(V^\perp,e_j)
=\overline g(\overline\nabla_iV^\perp,e_j)
+\overline g(V^\perp,\overline\nabla_ie_j)
\]
得到
\[
\overline g(\overline\nabla_iV^\perp,e_j)
=-\overline g(V^\perp,B_{ij}).
\]
收缩后
\[
\frac12\operatorname{tr}_g\dot g
=\operatorname{div}_M V^\top
-m\overline g(\vb{H},V^\perp).
\]
再用体积形式的一般变分公式
\[
\partial_t\operatorname{vol}_g
=\frac12\operatorname{tr}_g\dot g\,\operatorname{vol}_g
\]
并对散度项积分，即得到\eqnrefs{eq:math-dg-first-variation-submanifold-volume}。

一阶变分 $\delta\operatorname{Area}=-\int_M m\,\overline g(\vb{H},V^\perp)\dd\operatorname{vol}_g$ 说明面积泛函的 Euler--Lagrange 方程为 $\vb{H}=0$，即极小子流形方程：切向变分只产生散度（边界项），法向变分才通过平均曲率与面积变化耦合。

典型例：平面 $B_{ij}=0$、$\vb{H}=0$；半径 $R$ 的球面取外法向时 $\vb{H}=-(1/R)n\neq0$；悬链面 $\vb{H}=0$，即两圆环间肥皂膜的平衡形状。
\end{proof}

\begin{theorem}[欧氏 极小超曲面的 Simons 恒等式]
设 $M^m\subset\RR^{m+1}$ 为光滑极小超曲面，$A$ 为形状算符。则\[\Delta_{\mathrm{LB}}h_{ij}=-|A|^2h_{ij},\qquad\tfrac12\Delta_{\mathrm{LB}}|A|^2=|\nabla A|^2-|A|^4.\]
\end{theorem}

\begin{proof}
在 欧氏 超曲面 上，Codazzi 给出
\[
\nabla_k h_{ij}=\nabla_i h_{kj}.
\]
因此
\[
\Delta_{\mathrm{LB}}h_{ij}
=\nabla^k\nabla_kh_{ij}
=\nabla^k\nabla_i h_{kj}.
\]
交换两次协变导数时产生 Riemann 曲率 对两个切指标的作用；再用 Gauss 方程把所有内在 Riemann 张量 写成 $h*h$。另一方面，极小 条件
\[
\operatorname{tr}A=0
\]
使 $\nabla_i\nabla_j(\operatorname{tr}A)$ 项完全消失。剩余的所有三次项整理后恰为
\[
-|A|^2h_{ij},
\]
从而得到\eqnrefs{eq:math-dg-simons-hessian}。最后对
\[
|A|^2=h_{ij}h^{ij}
\]
应用乘积法则：
\[
\frac12\Delta_{\mathrm{LB}}|A|^2
=|\nabla A|^2
+\langle\Delta_{\mathrm{LB}}A,A\rangle
=|\nabla A|^2-|A|^4.
\]
所以 Simons 恒等式 的三项来源分别非常清楚：Codazzi 负责交换一个导数，Gauss 负责把交换子曲率变成 $A^2$，minimality 负责杀掉 平均-曲率 Hessian。

交换子计算可在一点的主曲率正交标架中完整核对。
记 $H_0=\operatorname{tr}A$，在该点 $h_{ij}=\kappa_i\delta_{ij}$。
由 Codazzi 与导数交换，
\[
\Delta_{\mathrm{LB}}h_{ij}
=\nabla_i\nabla_jH_0+
\sum_k([\nabla_k,\nabla_i]h)_{kj}.
\]
两个协变指标各贡献一项；由 Gauss 方程，它们分别给出
\[
\sum_k([\nabla_k,\nabla_i]h)_{kj}
=H_0(A^2)_{ij}-|A|^2h_{ij}.
\]
对角分量的逐项计算尤其直接：$K(e_i,e_k)=\kappa_i\kappa_k$，因此
\[
\begin{aligned}
\sum_k([\nabla_k,\nabla_i]h)_{ki}
&=\sum_kK(e_i,e_k)(\kappa_i-\kappa_k)\\
&=\kappa_i^2\sum_k\kappa_k-\kappa_i\sum_k\kappa_k^2\\
&=H_0\kappa_i^2-|A|^2\kappa_i.
\end{aligned}
\]
非对角项由同一 Gauss 分量计算为零；张量性使该表达与标架无关。
所以一般 欧氏 超曲面满足
\[
\Delta_{\mathrm{LB}}h_{ij}
=\nabla_i\nabla_jH_0+H_0(A^2)_{ij}-|A|^2h_{ij}.
\]
极小条件 $H_0=0$ 才消去前两项。
这里 $\operatorname{Ric}=H_0h-A^2$，且一般 $A^2\ne|A|^2I$；
例如主曲率为 $1,-1$ 时，$A^2=I$ 而 $|A|^2=2$。
因此不能用这个错误的矩阵恒等式来压缩三次曲率项。
\end{proof}

